% Part: turing-machines
% Chapter: machines-computations
% Section: variants

\documentclass[../../../include/open-logic-section]{subfiles}

\begin{document}

\olfileid{tur}{mac}{var}
\olsection{گونه‌های ماشین تورینگ}

در واقع، ماشین‌های تورینگ را به شیوه‌های ممکنِ بسیاری می‌توان تعریف
کرد و تعریف ما تنها یکی از آن‌هاست. تعریف ما از بعضی جهت‌ها آزادتر از
تعریف‌های دیگر است. ما الفباهای متناهیِ دلخواه را مجاز می‌دانیم؛ در
تعریفی محدودتر شاید فقط دو نماد نواریِ~$\TMstroke$ و~$\TMblank$ مجاز
باشند. ما به ماشین اجازه می‌دهیم هم‌زمان نمادی روی نوار بنویسد و حرکت
کند؛ در تعریف‌های دیگر ماشین یا می‌نویسد یا حرکت می‌کند. امکان نوشتن
بدون حرکت سر نوار را نیز مجاز می‌دانیم؛ برخی تعریف‌ها «دستور»~$\TMstay$
را کنار می‌گذارند. از جهت‌هایی دیگر، تعریف ما محدودتر است. فرض کردیم
نوار تنها در یک جهت نامتناهی است؛ تعریف‌های دیگری نوار را هم به چپ و
هم به راست نامتناهی می‌گیرند. حتی می‌توان هر تعداد نوار جداگانه یا
شبکه‌ای نامتناهی از خانه‌ها را مجاز دانست. ما مجموعه‌دستورهای ماشین
تورینگ را با تابع گذار نمایش می‌دهیم؛ تعریف‌های دیگری رابطهٔ گذار به
کار می‌برند که در هر وضعیت بیش از یک دستور ممکن به ماشین می‌دهد.

این آخرین آزادسازیِ تعریف به‌ویژه جالب است. در تعریف ما، هنگامی که
ماشین در حالت~$q$ نماد~$\sigma$ را می‌خواند، $\delta(q,\sigma)$ نماد
جدید، حالت جدید و مکان جدید سر نوار را تعیین می‌کند. اما اگر
مجموعه‌دستور را رابطه‌ای میان زوج‌های حالت--نماد فعلیِ
$\tuple{q,\sigma}$ و سه‌تایی‌های حالت--نماد--جهت جدیدِ
$\tuple{q',\sigma',D}$ بگیریم، ممکن است عمل ماشین تورینگ به‌طور یکتا
تعیین نشود: رابطهٔ دستور می‌تواند هم
$\tuple{q,\sigma,q',\sigma',D}$ و هم
$\tuple{q,\sigma,q'',\sigma'',D'}$ را در بر داشته باشد. در این صورت
یک ماشین تورینگِ \emph{نامعین} داریم. این ماشین‌ها در نظریهٔ پیچیدگی
محاسباتی نقشی مهم دارند.

دربارهٔ زمان توقف ماشین تورینگ نیز قراردادهای گوناگونی وجود دارد: ما
می‌گوییم ماشین هنگامی متوقف می‌شود که تابع گذار تعریف‌نشده باشد؛ در
تعریف‌های دیگر ماشین باید در یک حالت توقفِ ویژه و معین باشد. در
\olref[hal]{sec} توضیح دادیم که الزامِ حالت توقف معین، محدودیتی نیست
که بر آنچه ماشین‌های تورینگ می‌توانند محاسبه کنند اثر بگذارد. چون
نوارهای ماشین‌های تورینگ ما تنها در یک جهت نامتناهی‌اند، مواردی هست
که ماشین نمی‌تواند دستوری را به‌درستی اجرا کند: وقتی چپ‌ترین خانه را
می‌خواند و باید به چپ برود. بنا بر تعریف ما، به‌جای «افتادن از لبه»
همان‌جا می‌ماند؛ اما می‌توانستیم تعریف کنیم که در این وضعیت متوقف
شود. این تعریف نیز هم‌ارز است: می‌توانیم رفتار ماشینی را که هنگام
تلاش برای حرکت از خانهٔ~$0$ به چپ متوقف می‌شود با حذف هر گذار

\[
\delta(q,\TMendtape)=\tuple{q',\sigma,\TMleft}
\]

شبیه‌سازی کنیم؛ در این صورت ماشین به‌جای تلاش برای حرکت به چپ روی
$\TMendtape$ متوقف می‌شود.\footnote{این راه \emph{کاملاً} درست نیست،
زیرا چیزی مانع نوشتن و خواندن~$\TMendtape$ در خانه‌هایی غیر از
خانهٔ~$0$ نمی‌شود (\olref[una]{ex:mover} را ببینید). می‌توان با افزودن
نماد دومی چون~$\TMendtape'$ که برای چنین منظوری به کار رود، این مشکل
را برطرف کرد.}

عددها، و در نتیجه تابع ورودی--خروجیِ محاسبه‌شده به‌وسیلهٔ ماشین
تورینگ، را نیز به روش‌های گوناگون می‌توان نمایش داد: ما نمایش یکانی
را به کار می‌بریم، اما نمایش دودویی نیز ممکن است. این نمایش افزون بر
$\TMblank$ و~$\TMendtape$ به دو نماد نیاز دارد.

اکنون واقعیتی جالب در میان است: هیچ‌یک از این گونه‌ها در اینکه چه
توابعی تورینگ‌محاسبه‌پذیرند تفاوتی ایجاد نمی‌کند. \emph{اگر تابعی
طبق یک تعریف تورینگ‌محاسبه‌پذیر باشد، طبق همهٔ آن‌ها
تورینگ‌محاسبه‌پذیر است.}

به جزئیات اثبات این مطلب نمی‌پردازیم. فقط یک نمونه: با مجازکردن نواری
که در هر دو جهت نامتناهی است یا چند نوار، توان محاسباتی بیشتری به دست
نمی‌آوریم. دلیل، به‌طور تقریبی، این است که ماشین تورینگ دارای یک نوار
یک‌سویهٔ نامتناهی می‌تواند چند نوار یا نوار دوسویهٔ نامتناهی را
شبیه‌سازی کند. برای نمونه، با به‌کارگیری حالت‌ها و دستورهای اضافی
می‌توان برنامهٔ ماشینی با چند نوار یا نوار دوسویهٔ نامتناهی را به
برنامهٔ ماشینی با یک نوار یک‌سویهٔ نامتناهی «ترجمه» کرد. ماشین
ترجمه‌شده می‌تواند خانه‌های زوج را برای خانه‌های نوار~$1$ (یا
خانه‌های «مثبت» نوار دوسویهٔ نامتناهی) و خانه‌های فرد را برای
خانه‌های نوار~$2$ (یا خانه‌های «منفی») به کار ببرد.

\end{document}
